% Standard ACM conference proceedings: two-column sigconf layout.
% Compile with latexmk -pdf main.tex (BibTeX runs automatically).
\documentclass[sigconf]{acmart}

% Keep the class defaults for margins, typography, and bibliography.
\AtBeginDocument{%
  \providecommand\BibTeX{{Bib\TeX}}}

% DRAFT metadata. Replace with the exact ACM rights-form information
% before preparing the final publication version.
\setcopyright{none}
\acmYear{2026}
\copyrightyear{2026}
\acmConference[CONF '26]{Conference Name (Replace Before Submission)}%
  {Month DD--DD, 2026}{City, Country}
\acmBooktitle{Conference Name (CONF '26), Month DD--DD, 2026, City, Country}
\acmDOI{}
\acmISBN{}

\begin{document}

\title[Short Paper Title]{Your Paper Title: A Descriptive Subtitle}

\author{First Author}
\email{first.author@example.com}
\affiliation{%
  \institution{University or Organization}
  \city{City}
  \country{Country}
}

\author{Second Author}
\email{second.author@example.com}
\affiliation{%
  \institution{University or Organization}
  \city{City}
  \country{Country}
}

\renewcommand{\shortauthors}{First Author and Second Author}

\begin{abstract}
Replace this text with a self-contained summary of your research problem,
approach, main findings, and implications. State the problem and its importance,
describe the proposed method, summarize the evidence, and explain the principal
contribution. This file is an editable paper scaffold; all example content,
author details, and conference metadata must be replaced before submission.
\end{abstract}

% Replace these example CCS terms using https://dl.acm.org/ccs .
\begin{CCSXML}
<ccs2012>
 <concept>
  <concept_id>10010147</concept_id>
  <concept_desc>Computing methodologies</concept_desc>
  <concept_significance>500</concept_significance>
 </concept>
</ccs2012>
\end{CCSXML}
\ccsdesc[500]{Computing methodologies}
\keywords{keyword one, keyword two, keyword three}

\maketitle

\section{Introduction}
Introduce the research problem, explain why it matters, and identify the gap
addressed by this paper. Define the intended scope and explain the practical or
theoretical setting. Keep the claims specific enough to be supported by the
evaluation.

Summarize your approach and list the main contributions:
\begin{itemize}
  \item A clearly stated research question or problem formulation.
  \item A method, system, or analysis that addresses the identified gap.
  \item An evaluation that supports the paper's principal claims.
\end{itemize}

The remaining sections provide a starting structure for a conference paper.
Replace the instructional prose with your own research.

\section{Related Work}
Organize prior research by theme and explain how your contribution relates to
each theme. Use citations to support specific statements. For example,
\citet{lamport1994latex} provides a reference for document preparation in
\LaTeX. This citation demonstrates the ACM bibliography style; it is not
intended as evidence for your research topic.

\section{Method}
\subsection{Problem Formulation}
Define the inputs, outputs, assumptions, and notation. For example, a training
objective may be written as
\begin{equation}
  \mathcal{L}(\theta)
  = \frac{1}{N}\sum_{i=1}^{N}\ell(f_{\theta}(x_i),y_i),
  \label{eq:objective}
\end{equation}
where $N$ is the number of examples and $\theta$ denotes model parameters.
Replace Equation~\ref{eq:objective} with an expression appropriate to your work.

\subsection{Proposed Approach}
Describe the core idea and the steps needed to reproduce it. Explain design
choices and provide sufficient detail about implementation or analysis.

\begin{figure}[t]
  \centering
  \fbox{\parbox[c][22mm][c]{0.90\columnwidth}{\centering
    Replace with your figure\\
    Input $\longrightarrow$ Method $\longrightarrow$ Output}}
  \caption{A placeholder for a single-column overview figure. Replace the box
    with \texttt{\string\includegraphics[width=\string\linewidth]\{images/figure.pdf\}}.}
  \Description{A placeholder diagram showing input, method, and output.}
  \label{fig:overview}
\end{figure}

Figure~\ref{fig:overview} illustrates the approach. Give every figure a useful
caption and an accessibility description.

\section{Evaluation}
\subsection{Experimental Setup}
Specify datasets, baselines, metrics, hardware, and relevant parameter choices.
Describe the procedure and any statistical analysis. Explain how results can
be reproduced.

\subsection{Results}
Report observations supported by the collected evidence.
Table~\ref{tab:results} provides a formatting example. Its entries are
placeholders, not experimental findings.

\begin{table}[t]
  \caption{Example results structure. Replace all placeholder entries.}
  \label{tab:results}
  \centering
  \begin{tabular}{lcc}
    \toprule
    Method & Metric A & Metric B \\
    \midrule
    Baseline & --- & --- \\
    Proposed method & --- & --- \\
    \bottomrule
  \end{tabular}
\end{table}

\section{Discussion and Limitations}
Interpret the results in relation to the research question. Discuss
assumptions, threats to validity, failure cases, and the limits of
generalization. Include ethical considerations where relevant.

\section{Conclusion}
Summarize the contribution and the evidence that supports it. State the main
implication and identify a concrete direction for future work.

\begin{acks}
Replace this text with funding and acknowledgments, or remove this environment.
Omit identifying acknowledgments when preparing an anonymous submission.
\end{acks}

\bibliographystyle{ACM-Reference-Format}
\bibliography{references}

\end{document}
